\documentclass{article}
\usepackage{graphicx}
\usepackage{amsmath}
\usepackage{array}
\usepackage{fancyhdr}
\usepackage{amssymb}
\usepackage[shortlabels]{enumitem}

\pagestyle{fancy}
\fancyhead[L]{Banghao Chi}
\fancyhead[C]{Homework 5}
\fancyhead[R]{5th Mar}

\fancyfoot[C]{\thepage}

\renewcommand{\headrulewidth}{0.5pt}
\renewcommand{\footrulewidth}{0.5pt}

\begin{document}

\section*{Exercise 1}
Let $(x_n)_{n \in \mathbb{N}}$ and $(y_n)_{n \in \mathbb{N}}$ be two sequences of real numbers. Define the sequence $(z_n)_{n \in \mathbb{N}}$ by
\begin{align*}
z_{2n-1} &:= x_n\\
\end{align*}
and
\begin{align*}
z_{2n} &:= y_n
\end{align*}
for all $n \in \mathbb{N}$. Show that $(z_n)_{n \in \mathbb{N}}$ is convergent if and only if both $(x_n)_{n \in \mathbb{N}}$ and $(y_n)_{n \in \mathbb{N}}$ are convergent and
\begin{align*}
\lim x_n = \lim y_n.
\end{align*}

\textbf{Solution: }\\

Since this is a if and only if statement, we will show both directions. \\

($\Rightarrow$) $(z_{2n-1})_{n \in \mathbb{N}}$ and $(z_{2n})_{n \in \mathbb{N}}$ are subsequences of $(z_n)_{n \in \mathbb{N}}$. Since $(z_n)_{n \in \mathbb{N}}$ is convergent, all of its subsequences must converge to the same limit $L$. Therefore:
\begin{align*}
\lim z_{2n-1} = L \quad \text{and} \quad \lim z_{2n} = L
\end{align*}

Since $z_{2n-1} = x_n$ and $z_{2n} = y_n$ for all $n \in \mathbb{N}$, we have:
\begin{align*}
\lim x_n = L \quad \text{and} \quad \lim y_n = L
\end{align*}

which proves that both sequences $(x_n)_{n \in \mathbb{N}}$ and $(y_n)_{n \in \mathbb{N}}$ are convergent and $\lim x_n = \lim y_n = L$. \\

($\Leftarrow$) Suppose both $(x_n)_{n \in \mathbb{N}}$ and $(y_n)_{n \in \mathbb{N}}$ are convergent with $\lim x_n = \lim y_n = L$. \\

Let $\epsilon > 0$. Since $\lim x_n = L$, there exists $N_1 \in \mathbb{N}$ such that for all $n \geq N_1$:
\begin{align*}
|x_n - L| < \epsilon
\end{align*}

Similarly, since $\lim y_n = L$, there exists $N_2 \in \mathbb{N}$ such that for all $n \geq N_2$:
\begin{align*}
|y_n - L| < \epsilon
\end{align*}

Let $N = \max(2N_1 - 1, 2N_2)$. Then for all $n \geq N$: \\

If $n = 2k-1$, then $k = \frac{n+1}{2} \geq \frac{2N_1-1+1}{2} = N_1$. Thus:
\begin{align*}
|z_n - L| = |z_{2k-1} - L| = |x_k - L| < \epsilon
\end{align*}

If $n = 2k$, then $k = \frac{n}{2} \geq \frac{2N_2}{2} = N_2$. Thus:
\begin{align*}
|z_n - L| = |z_{2k} - L| = |y_k - L| < \epsilon
\end{align*}

Therefore, $\lim z_n = L$, which proves that $(z_n)_{n \in \mathbb{N}}$ is convergent. \\

Hence, $(z_n)_{n \in \mathbb{N}}$ is convergent if and only if both $(x_n)_{n \in \mathbb{N}}$ and $(y_n)_{n \in \mathbb{N}}$ are convergent and $\lim x_n = \lim y_n$.

\newpage

\section*{Exercise 2}
Show that if $(x_n)_{n \in \mathbb{N}}$ is not bounded, then there exists a subsequence $(y_n)_{n \in \mathbb{N}}$ of $(x_n)_{n \in \mathbb{N}}$ such that
\begin{align*}
\lim \frac{1}{y_n} = 0.
\end{align*}

\textbf{Solution: }\\

Construct a subsequence $(y_n)_{n \in \mathbb{N}}$ of $(x_n)_{n \in \mathbb{N}}$ as follows: \\

Let $n_1$ be the smallest index such that $|x_{n_1}| > 1$. Existence holds because $(x_n)_{n \in \mathbb{N}}$ is unbounded. Let $y_1 = x_{n_1}$. \\

For $k \geq 2$, let $n_k$ be the smallest index with $n_k > n_{k-1}$ such that $|x_{n_k}| > k$. Existence holds because $(x_n)_{n \in \mathbb{N}}$ is unbounded. Let $y_k = x_{n_k}$. \\

By construction, $(y_k)_{k \in \mathbb{N}}$ is a subsequence of $(x_n)_{n \in \mathbb{N}}$ where $|y_k| > k$ for all $k \in \mathbb{N}$. \\

For all $k \in \mathbb{N}$, we have 
\begin{align*}
\left|\frac{1}{y_k}\right| = \frac{1}{|y_k|} < \frac{1}{k}
\end{align*}

Since $\lim \frac{1}{k} = 0$ and $\left|\frac{1}{y_k}\right| \geq 0$ for all $k \in \mathbb{N}$, by the squeeze theorem, 
\begin{align*}
\lim \left|\frac{1}{y_k}\right| = 0
\end{align*}

which implies that $\lim \frac{1}{y_k} = 0$, completing the proof.

\newpage

\section*{Exercise 3}
Define $(x_n)_{n \in \mathbb{N}}$ by
\begin{align*}
x_{2n-1} &:= 1+\frac{1}{n}\\
\end{align*}
and
\begin{align*}
x_{2n} &:= -\frac{1}{n}
\end{align*}
for all $n \in \mathbb{N}$. Compute
\begin{enumerate}[(a)]
\item $\lim \sup x_n$,
\item $\lim \inf x_n$,
\item $\sup\{x_n : n \in \mathbb{N}\}$, and
\item $\inf\{x_n : n \in \mathbb{N}\}$.
\end{enumerate}

\textbf{Solution: }\\

$(x_n)_{n \in \mathbb{N}}$ has two subsequences:
\begin{itemize}
    \item For odd terms: $\lim x_{2n-1} = \lim (1+\frac{1}{n}) = 1$
    \item For even terms: $\lim x_{2n} = \lim (-\frac{1}{n}) = 0$
\end{itemize}

(a) By theorem, $\lim \sup x_n$ is the larger of these two subsequential limits, so:
\begin{align*}
\lim \sup x_n = 1
\end{align*}

(b) By theorem, $\lim \inf x_n$ is the smaller of the two subsequential limits, so:
\begin{align*}
\lim \inf x_n = 0
\end{align*}

(c) The supremum of the entire sequence is the largest term of the sequence. Since even-indexed term is negative, we only need to consider the odd-indexed terms. For the odd-indexed terms $x_{2n-1} = 1+\frac{1}{n}$, this value is largest when $n=1$, giving:
\begin{align*}
\sup\{x_n : n \in \mathbb{N}\} = x_1 = 2
\end{align*}

(d) The infimum of the entire sequence is the smallest term of the sequence. Since odd-indexed term is positive, we only need to consider the even-indexed terms. For the even-indexed terms $x_{2n} = -\frac{1}{n}$, this value is smallest when $n=1$, giving:
\begin{align*}
\inf\{x_n : n \in \mathbb{N}\} = x_2 = -1
\end{align*}

\newpage

\section*{Exercise 4}
If $0 < r < 1$ and
\begin{align*}
|x_{n+1} - x_n| < r^n
\end{align*}
for all $n \in \mathbb{N}$, show that $(x_n)_{n \in \mathbb{N}}$ is a Cauchy sequence. \\

\textbf{Solution: }\\

We want to show that for any $\epsilon > 0$, there exists an $N \in \mathbb{N}$ such that for all $m, n \geq N$, we have $|x_m - x_n| < \epsilon$. \\

Without loss of generality, assume $m > n$ (since if $m = n$, then $|x_m - x_n| = 0 < \epsilon$). \\

For any $m, n \in \mathbb{N}$ with $m > n$, we have:
\begin{align*}
|x_m - x_n| &= |x_m - x_{m-1} + x_{m-1} - x_{m-2} + \cdots + x_{n+1} - x_n| \\
&\leq |x_m - x_{m-1}| + |x_{m-1} - x_{m-2}| + \cdots + |x_{n+1} - x_n| \\
&< r^{m-1} + r^{m-2} + \cdots + r^n \\
&= r^n(1 + r + r^2 + \cdots + r^{m-n-1}) \\
&= r^n \frac{1 - r^{m-n}}{1 - r} \\
&= \frac{r^n - r^m}{1 - r} \\
&< \frac{r^n}{1 - r}
\end{align*}

Here, since $0 < r < 1$, we have $1 - r > 0$, and $0 < r^n < 1$ for all $n \in \mathbb{N}$. \\

And since $n$ is arbitrary, for any $\epsilon > 0$, we always choose $N \in \mathbb{N}$ such that:
\begin{align*}
\frac{r^N}{1 - r} < \epsilon
\end{align*}

Then for all $m, n \geq N$, we have:
\begin{align*}
|x_m - x_n| < \frac{r^N}{1 - r} < \epsilon
\end{align*}

Therefore, $(x_n)_{n \in \mathbb{N}}$ is a Cauchy sequence.

\newpage

\section*{Exercise 5}
Define $(x_n)_{n \in \mathbb{N}}$ by
\begin{align*}
x_1 &:= 2\\
\end{align*}
and
\begin{align*}
x_{n+1} &:= 2 + \frac{1}{x_n}
\end{align*}
for all $n \in \mathbb{N}$. Show that $(x_n)_{n \in \mathbb{N}}$ converges and find the limit. \\

\textbf{Solution: }
\begin{align*}
x_{n+1} - x_n &= 2 + \frac{1}{x_n} - x_n\\
&= \frac{2x_n + 1 - x_n^2}{x_n}\\
&= \frac{-(x_n^2 - 2x_n - 1)}{x_n}
\end{align*}

The quadratic expression $x_n^2 - 2x_n - 1$ has roots at $1 \pm \sqrt{2}$. Since $x_n > 0$ for all $n$, we only need to consider the positive root $1 + \sqrt{2}$. \\

For the sign of $x_{n+1} - x_n$:
\begin{align*}
\text{If } 0 < x_n < 1 + \sqrt{2} &\implies x_n^2 - 2x_n - 1 < 0 \implies x_{n+1} - x_n > 0\\
\text{If } x_n > 1 + \sqrt{2} &\implies x_n^2 - 2x_n - 1 > 0 \implies x_{n+1} - x_n < 0\\
\text{If } x_n = 1 + \sqrt{2} &\implies x_n^2 - 2x_n - 1 = 0 \implies x_{n+1} - x_n = 0
\end{align*}

Simplifying the above, we get:
\begin{align*}
\text{If } x_n < 1 + \sqrt{2} &\implies x_n < x_{n+1}\\
\text{If } x_n > 1 + \sqrt{2} &\implies x_n > x_{n+1}\\
\text{If } x_n = 1 + \sqrt{2} &\implies x_n = x_{n+1}
\end{align*}

Computing the first few terms:
\begin{align*}
x_1 &= 2 < 1 + \sqrt{2} \approx 2.414 \implies x_1 < x_2\\
x_2 &= 2.5 > 1 + \sqrt{2} \implies x_2 > x_3\\
x_3 &= 2.4 < 1 + \sqrt{2} \implies x_3 < x_4\\
x_4 &\approx 2.417 > 1 + \sqrt{2} \implies x_4 > x_5
\end{align*}

Observing the sequence, we see that it oscillates around $1 + \sqrt{2}$ with:
\begin{align*}
x_1 < x_3 < x_5 < \ldots < 1 + \sqrt{2} < \ldots < x_6 < x_4 < x_2
\end{align*}

Therefore, we want to show that these two subsequences(odd-indexed and even-indexed subsequences, which form the entire sequence) converge to the same limit $1 + \sqrt{2}$. \\

Prove by induction that for all $n \geq 1$:
\begin{itemize}
    \item $x_{2n-1} < 1 + \sqrt{2}$ (all odd-indexed terms are less than $1 + \sqrt{2}$)
    \item $x_{2n} > 1 + \sqrt{2}$ (all even-indexed terms are greater than $1 + \sqrt{2}$)
\end{itemize}

Base case: $x_1 = 2 < 1 + \sqrt{2}$ and $x_2 = 2.5 > 1 + \sqrt{2}$. \\

Inductive step: Assume that for some $k \geq 1$, we have:
\begin{align*}
x_{2k-1} < 1 + \sqrt{2} \\
x_{2k} > 1 + \sqrt{2}
\end{align*}

For $x_{2k+1}$:
\begin{align*}
x_{2k+1} = 2 + \frac{1}{x_{2k}}
\end{align*}

Since $x_{2k} > 1 + \sqrt{2}$, we have $\frac{1}{x_{2k}} < \frac{1}{1 + \sqrt{2}} = \sqrt{2} - 1$. \\

Therefore, $x_{2k+1} = 2 + \frac{1}{x_{2k}} < 2 + \sqrt{2} - 1 = 1 + \sqrt{2}$. \\

For $x_{2k+2}$:
\begin{align*}
x_{2k+2} = 2 + \frac{1}{x_{2k+1}}
\end{align*}

Since $x_{2k+1} < 1 + \sqrt{2}$, we have $\frac{1}{x_{2k+1}} > \frac{1}{1 + \sqrt{2}} = \sqrt{2} - 1$. \\

Therefore, $x_{2k+2} = 2 + \frac{1}{x_{2k+1}} > 2 + \sqrt{2} - 1 = 1 + \sqrt{2}$. \\

By the principle of mathematical induction, the pattern holds for all $n \geq 1$. \\

Next, we use induction to prove both monotonicity properties. Define $f(x) = 2 + \frac{1}{x}$. \\

Base case: We've calculated $x_1 = 2$, $x_2 = 2.5$, $x_3 = 2.4$, and $x_4 \approx 2.417$.
We can verify that $x_3 > x_1$ and $x_4 < x_2$. \\

Inductive hypothesis: Assume for some $k \geq 2$ that $x_{2k-1} > x_{2k-3}$ and $x_{2k} < x_{2k-2}$. \\

For odd terms: We have $x_{2k+1} = f(x_{2k})$ and $x_{2k-1} = f(x_{2k-2})$.
By our inductive hypothesis, $x_{2k} < x_{2k-2}$.
Since $f$ is strictly decreasing, $f(x_{2k}) > f(x_{2k-2})$, which means $x_{2k+1} > x_{2k-1}$. \\

For even terms: We have $x_{2k+2} = f(x_{2k+1})$ and $x_{2k} = f(x_{2k-1})$.
From what we just proved, $x_{2k+1} > x_{2k-1}$.
Since $f$ is strictly decreasing, $f(x_{2k+1}) < f(x_{2k-1})$, which means $x_{2k+2} < x_{2k}$. \\

By the principle of mathematical induction, both properties hold for all $n \geq 1$. \\

And by the Monotone Convergence Theorem, both subsequences converge. Denote their limits as $L_{odd}$ and $L_{even}$ respectively. \\

Using the recurrence relation $x_{n+1} = 2 + \frac{1}{x_n}$, we get:
\begin{align*}
L_{even} = 2 + \frac{1}{L_{odd}} \\
L_{odd} = 2 + \frac{1}{L_{even}}
\end{align*}

Substituting the first equation into the second:
\begin{align*}
L_{odd} = 2 + \frac{1}{2 + \frac{1}{L_{odd}}}
\end{align*}

We end up with a quadratic equation:
\begin{align*}
L_{odd}^2 - 2L_{odd} - 1 = 0
\end{align*}

Using the quadratic formula:
\begin{align*}
L_{odd} = \frac{2 \pm \sqrt{4 + 4}}{2} = 1 \pm \sqrt{2}
\end{align*}

Since $L_{odd} > 0$ (all terms in our sequence are positive), we have $L_{odd} = 1 + \sqrt{2}$. \\

Similarly, we solve for $L_{even}$ and obtain $L_{even} = 1 + \sqrt{2}$. \\

Since both subsequences converge to the same limit $1 + \sqrt{2}$, and they contain all terms of the original sequence, the sequence $(x_n)_{n \in \mathbb{N}}$ converges to $1 + \sqrt{2}$.

\end{document}